\documentclass[10pt,a4paper]{amsart}
\usepackage[margin=19mm]{geometry}
\usepackage{amsmath,amssymb,mathtools}
\usepackage[scr=rsfs,cal=euler]{mathalfa}
\usepackage{xfrac}
\usepackage[unicode,hidelinks,bookmarks=false]{hyperref}
\newcommand{\ie}{i.e.\ }
\newcommand{\cf}{cf.\ }
\newcommand{\resp}{respectively\ }
\newcommand{\Subsection}{\subsection}
\providecommand{\nobreakdash}{\nobreak\hbox{-}}
\setlength{\emergencystretch}{2em}
\multlinegap0pt
\usepackage{amssymb}
\usepackage{tikz}
\usepackage{dsfont}
\newtheorem{theorem}{Theorem}
\newcommand{\EtO}{\mathrm{EtO}}
\newcommand{\OtE}{\mathrm{OtE}}
\newcommand{\QSym}{\mathrm{QSym}}
\title{Extending the descent-to-peak map and its applications}
\author{Farid Aliniaeifard, Shu Xiao Li}
\date{Source: arXiv:2110.05648v2 (CC BY 4.0)}
\begin{document}
\maketitle

\noindent\emph{Source and licence.} Extending the descent-to-peak map and its applications. Version arXiv:2110.05648v2 (CC BY 4.0), distributed by arXiv under the Creative Commons Attribution 4.0 International licence, http://creativecommons.org/licenses/by/4.0/, as recorded in the arXiv licence metadata for this exact version and shown on the abstract page https://arxiv.org/abs/2110.05648v2. Full licence notice: \texttt{licenses/arxiv-2110.05648-CC-BY-4.0.txt}.\par\smallskip

\noindent\textit{Editorial scope.} Counted unit: the article's theorem thm:theta (source0.tex lines 630--636). The article's complete original proof of it is delivered with it (source0.tex lines 638--693, one \texttt{proof} environment of 56 lines). No mathematics is written, repaired or re-derived: the statement and the proof are the article's own lines, in the article's own notation, and no retained line is substituted or re-flowed.  No further source-local context is needed: every cross-reference of every retained line resolves inside the two delivered spans.  Nothing was omitted from any retained span, so the package carries no omission record.  No retained line contains a citation, so the wrapper carries no bibliography and no bibliography entry.  No retained line contains a figure environment, an image-inclusion command or drawing code, so no image is delivered and the package declares no asset dependency.  Editorial formatting only, in addition: an AMS \texttt{amsart} preamble that restates the article's own declarations and macros used by the retained lines, namely the environments \texttt{theorem} on separate counters, together with the standard packages those lines need.  The retained environments print in this document as Theorem 1 in order of appearance, whereas the article prints the same items with the numbering of its own full document.  The complete native source of the article as distributed in the arXiv e-print for this exact version is bundled unchanged as \texttt{sources/117075/source0.tex}.
\begin{theorem}\label{thm:theta}
	The map $\Theta_\QSym$ acts on the shuffle basis as follows
	$$\Theta_\QSym(S_\alpha)=\left\{\begin{array}{cc}
		2^{\ell(\alpha)}S_\alpha & \text{if }\alpha \text{ is odd}\\
		0& \text{otherwise.}
		\end{array}\right.$$
\end{theorem}

\begin{proof}
	Let $\Phi:\QSym\to\QSym$ be the linear map defined on the shuffle basis as
	$$\Phi(S_\alpha)=\left\{\begin{array}{cc}
	2^{\ell(\alpha)}S_\alpha & \text{if }\alpha \text{ is odd}\\
	0& \text{otherwise.}
	\end{array}\right.$$
	It is easy to check that $\Phi$ is a Hopf algebra morphism.
	Since $\Theta_\QSym$ is the unique combinatorial Hopf morphism from $(\QSym,\overline{\zeta_\QSym^{-1}}\zeta_\QSym)$ to $(\QSym,\zeta_\QSym)$, it suffices to show that $\overline{\zeta_\QSym^{-1}}\zeta_\QSym=\zeta_\QSym\circ\Phi$.
	
	Recall that
	$$\zeta_\QSym(M_\alpha)=\left\{\begin{array}{cc}
	1 & \text{if }\ell(\alpha)=0\text{ or }1\\
	0 & \text{otherwise}.
	\end{array}\right.$$
	It follows from the definition of shuffle functions that,
	\begin{equation}\label{eq:3.1}
		\zeta_\QSym(S_\alpha)=\left\{\begin{array}{cc}
			\cfrac{1}{p!(\ell(\alpha)-p)!} & \text{if }\EtO(\alpha)=\{p\}\text{ and }\OtE(\alpha)=\emptyset\\
			\cfrac{1}{\ell(\alpha)!} & \text{if }\EtO(\alpha)=\OtE(\alpha)=\emptyset\\
			0 & \text{otherwise}.
		\end{array}\right.
	\end{equation}
	Hence, we obtain the following formula for $\zeta_\QSym\circ\Phi$
	$$\zeta_\QSym\circ\Phi(S_\alpha)=\left\{\begin{array}{cc}
	\cfrac{2^{\ell(\alpha)}}{\ell(\alpha)!}& \text{if }\alpha\text{ is odd}\\
	0 & \text{otherwise.}
	\end{array}\right.$$
	Combining with the antipode formula yields
	\begin{equation}\label{eq:3.2}
		\overline{\zeta_\QSym^{-1}}(S_\alpha)=\left\{\begin{array}{cc}
			(-1)^{|\alpha|+\ell(\alpha)}\cfrac{1}{q!(\ell(\alpha)-q)!} & \text{if }\EtO(\alpha)=\emptyset\text{ and }\OtE(\alpha)=\{q\}\\
			(-1)^{|\alpha|+\ell(\alpha)}\cfrac{1}{\ell(\alpha)!} & \text{if }\EtO(\alpha)=\OtE(\alpha)=\emptyset\\
			0 & \text{otherwise}.
		\end{array}\right.
	\end{equation}
	To sum up,
	\begin{align*}
	&\overline{\zeta_\QSym^{-1}}\zeta_\QSym(S_\alpha)=\sum_{i=0}^{\ell(\alpha)}\overline{\zeta_\QSym^{-1}}(S_{(\alpha_1,\dots,\alpha_i)})\zeta_\QSym(S_{(\alpha_{i+1},\dots,\alpha_{\ell(\alpha)})})\\
	=&\left\{\begin{array}{cc}
	\displaystyle\sum_{i=q}^p(-1)^{\alpha_1+\cdots+\alpha_i+i}\frac{1}{q!(i-q)!}\frac{1}{(p-i)!(\ell(\alpha)-p)!} & \text{if }\EtO(\alpha)=\{p\},\OtE(\alpha)=\{q\}\text{ and }q<p\\
	\displaystyle\sum_{i=0}^p(-1)^{\alpha_1+\cdots+\alpha_i+i}\frac{1}{i!}\frac{1}{(p-i)!(\ell(\alpha)-p)!} & \text{if }\EtO(\alpha)=\{p\}\text{ and }\OtE(\alpha)=\emptyset\\
	\displaystyle\sum_{i=q}^{\ell(\alpha)}(-1)^{\alpha_1+\cdots+\alpha_i+i}\frac{1}{q!(i-q)!}\frac{1}{(\ell(\alpha)-i)!} & \text{if }\EtO(\alpha)=\emptyset\text{ and }\OtE(\alpha)=\{q\}\\
	\displaystyle\sum_{i=0}^{\ell(\alpha)}(-1)^{\alpha_1+\cdots+\alpha_i+i}\frac{1}{i!}\frac{1}{(\ell(\alpha)-i)!} & \text{if }\EtO(\alpha)=\OtE(\alpha)=\emptyset\\
	0 & \text{otherwise.}
	\end{array}\right.\\
	\end{align*}
	In the first case, the summation is alternating in sign because $\alpha_i$ are even for $q< i\leq p$. Therefore, the summation can be simplified to
	$$\frac{(-1)^q}{q!(\ell(\alpha)-p)!}\sum_{i=q}^p(-1)^{i}\frac{1}{(i-q)!(p-i)!}.$$
	It is easy to see that this summation must vanish, e.g., multiply it by $(p-q)!$ and use the binomial identity.
	
	The same argument can be applied to all other cases except the sub-case of 4 that $\alpha$ is an odd composition. When $\alpha$ is odd, all terms in the sum have a positive sign. Hence, we have
	
	$$\overline{\zeta_\QSym^{-1}}\zeta_\QSym(S_\alpha)=\sum_{i=0}^{\ell(\alpha)}\frac{1}{i!(\ell(\alpha)-i)!}=\frac{2^{\ell(\alpha)}}{\ell(\alpha)!}.$$
	
	Comparing with the formula of $\zeta_\QSym\circ\Phi$ completes the proof.
\end{proof}

\end{document}
